\chapter{Introducción}
\label{cap:introduccion}

La seguridad de las aplicaciones web sigue siendo una de
las superficies de ataque más explotadas en incidentes de seguridad reales.
Informes como el OWASP Top~10~\citelib{OWASP21} sitúan de forma recurrente
vulnerabilidades como la inyección de código, el scripting o el
control de acceso roto entre los riesgos más críticos y más frecuentes en
aplicaciones, a pesar de tratarse de fallos ampliamente
documentados y, mayormente, los más sencillos de corregir una vez
identificados. Esta paradoja es el punto de partida de este trabajo.

Frente a un enfoque puramente teórico, este TFM
se acerca a estas vulnerabilidades desde la práctica: diseñando una
aplicación web propia, \breachlab, en la que cada fallo se presenta en la misma
de forma deliberada, se explota siguiendo una metodología de seguridad ofensiva
y se corrige aplicando las medidas estándar del sector, de modo que el
mismo entorno sirve para mostrar tanto el problema como la solución.

Además del componente práctico, este trabajo parte del estudio de una
metodología de referencia ya consolidada en el sector, la OWASP Web
Security Testing Guide (WSTG)~\citelink{OWASPWSTG}, para derivar de ella una
metodología propia más ágil: un subconjunto acotado de pruebas, centrado en
las categorías de mayor impacto, pensado para escenarios con recursos
limitados (equipos pequeños) en los que ejecutar la guía completa no
resulta viable. Dicha metodología propia se aplica sobre \breachlab de forma
controlada y se compara, en cobertura y en esfuerzo estimado, con la
metodología de referencia de la que parte (capítulos~\ref{cap:metodologia}
y~\ref{cap:comparativa}).

\section{Motivación}

La motivación de este trabajo es doble. Por un lado, personal y formativa:
comprender en profundidad cómo se manifiestan en código real vulnerabilidades
estudiadas comúnmente de forma conceptual: \emph{qué falla}, \emph{por qué
falla} y \emph{cómo se corrige}. Por otro lado,
práctica: la mayoría de entornos de entrenamiento disponibles públicamente
(DVWA, OWASP Juice Shop, entre otros) se apoyan en arquitecturas monolíticas
clásicas, mientras que buena parte de las aplicaciones actuales se despliegan
como una interfaz de tipo SPA (\emph{Single Page Application}) que aprovecha por ello
una API individual e independiente. Esta diferencia arquitectónica no es trivial:
modifica, por ejemplo, dónde se materializa un XSS o cómo se comporta la cookie de sesión
frente a un ataque CSRF.

Construir un laboratorio propio con esta arquitectura permite, además,
instrumentar el propio código para generar evidencias objetivas de cada
explotación mediante un registro de eventos, pudiendo alternar entre la rama
vulnerable y la rama corregida de cada endpoint sin la necesidad de mantener 
dos aplicaciones distintas, lo que resulta especialmente útil de cara a la 
documentación de este trabajo y posibles evoluciones futuras.

\section{Objetivos}
\label{sec:objetivos}

De acuerdo con la propuesta de este Trabajo Fin de Máster, el objetivo
general es estudiar una metodología de referencia para pruebas de seguridad
web, derivar de ella una metodología propia más ágil y validar ambas de
forma práctica mediante la creación de un entorno de pruebas controlado.
Dicho objetivo se materializaría
en los siguientes objetivos:

\begin{itemize}
\item Estudiar una metodología de referencia existente en el ámbito de la
  seguridad de aplicaciones web (OWASP WSTG).
\item Diseñar un entorno web vulnerable controlado.
\item Derivar, a partir de la metodología estudiada, una metodología propia
  orientada a la agilidad, acotando su alcance a un subconjunto de pruebas
  relevante en contextos con recursos limitados.
\item Identificar vulnerabilidades siguiendo dicha metodología propia.
\item Explotar las vulnerabilidades de forma controlada.
\item Analizar el impacto técnico y funcional de cada vulnerabilidad.
\item Comparar los resultados y el esfuerzo de la metodología propia frente
  a la metodología de referencia de la que parte.
\end{itemize}

\section{Alcance y limitaciones}

Este trabajo se limita a la capa de aplicación web, lo que significa que no se abordan
vulnerabilidades de red, de infraestructura ni de sistema operativo. De las
diez categorías del OWASP Top 10 \citelib{OWASP21}, se estudiarán en profundidad las seis
que pueden reproducirse y demostrarse de forma autocontenida en
\breachlab~(inyección SQL, XSS, CSRF, fallos de autenticación, control de
acceso roto y configuración de seguridad insegura); resultando en que el resto de 
vulnerabilidades sean comentadas de forma anecdótica únicamente como líneas de trabajo 
futuro (apartado~\ref{sec:trabajo-futuro}).

Por otro lado, todas las puntuaciones CVSS v3.1 calculadas a lo largo de la memoria son
ilustrativas de un escenario de laboratorio, con datos y usuarios ficticios y
sin exposición a Internet, lo que implica que no constituyen una auditoría de un sistema 
en producción y no deben interpretarse como tal. Del mismo modo, la comparativa
del capítulo~\ref{cap:comparativa} con DVWA y OWASP Juice Shop se restringe a
las vulnerabilidades tratadas en \breachlab, y no pretende ser una evaluación
exhaustiva de esos dos proyectos.

\needspace{15\baselineskip}
\section{Estructura de la memoria}

La memoria se organiza en los siguientes capítulos:

\begin{itemize}
\item El capítulo~\ref{cap:marco-teorico} recoge el estado del arte: el
  catálogo OWASP Top~10, la OWASP Web Security Testing Guide y otras
  metodologías de seguridad ofensiva, y el estándar CVSS v3.1 utilizado para
  valorar la severidad de cada vulnerabilidad.
\item El capítulo~\ref{cap:planificacion} recoge la planificación seguida
  para desarrollar el trabajo, con sus fases y duración estimada.
\item El capítulo~\ref{cap:diseno} describe el diseño y la arquitectura del
  entorno de pruebas \breachlab.
\item El capítulo~\ref{cap:metodologia} detalla la derivación de una
  metodología propia y ágil a partir de la OWASP WSTG, así como el proceso
  y las herramientas empleadas para aplicarla.
\item El capítulo~\ref{cap:explotacion} constituye el núcleo técnico del
  trabajo: la explotación y el análisis de impacto de cada una de las seis
  vulnerabilidades introducidas en \breachlab, siguiendo la metodología
  propia del capítulo~\ref{cap:metodologia}.
\item El capítulo~\ref{cap:comparativa} compara la metodología propia con la
  OWASP WSTG completa de la que parte, y compara además los resultados
  obtenidos con dos laboratorios de referencia, DVWA y OWASP Juice Shop.
\item El capítulo~\ref{cap:caso-estudio} presenta un caso de estudio que
  encadena varias de las vulnerabilidades anteriores en un escenario de
  ataque realista.
\item El capítulo~\ref{cap:conclusiones} recoge las conclusiones del trabajo
  y propone líneas de ampliación aplicables para futuras actualizaciones.
\end{itemize}
